\documentclass{gescons}

\genre {Entrevista}
\author{Gunther Kissmann}
\title{Diário de um Iniciante na Tenepes}


\begin{document}
    \makeentrevistatitle
    \coverart{back/Gunther_Kissmann}

    \begin{multicols}{2}

\begin{center}
    \includegraphics[width=8cm]{articles/entrevista/mockups/Gunther_Kissmann.png}
\end{center}


\textbf{1. Qual foi a motivação para a escrita da obra? Por que a definição deste tema para publicação de um livro?}

Tinha vontade de escrever um livro. Porém, tinha dúvidas sobre seu tema. Estimava que para a escrita, ter vivência sobre o assunto alicerçaria com mais consistência o texto em desenvolvimento.

Nos primeiros dias da prática da tenepes, percebi sugestão do amparador de função para organizar as anotações do diário de tenepes, de modo a ter substrato para escrita de livro.

A proposta desta obra é desdramatizar a prática da tenepes, descortinando a dinâmica ocorrida nos 60 primeiros dias de um iniciante na tenepes.

\begin{pullquote}
``A proposta desta obra é desdramatizar a prática da tenepes.''
\end{pullquote}

\textbf{2. Quais foram as principais percepções, intra e extrafísicas, durante a pesquisa e a escrita da obra? E posterior ao lançamento?}

Organizar um ambiente favorável à escrita cotidiana foi um aporte importante ao desenvolvimento do texto.

Em paralelo, foi perceptível em certos momentos o apoio do amparo extrafísico por meio de inspirações quanto aos caminhos ou ideias a serem acrescentadas no livro.

Notei pouco contrafluxo externo no processo de escrita, e houve períodos de insegurança sobre a relevância das informações que eu havia inserido no texto.

Após o lançamento, percebi momentos de evocação à minha pessoa, até o momento sempre de caráter positivo.

\textbf{3. Qual o maior aprendizado com a escrita desta obra?}

Sendo meu primeiro livro publicado, passei a ter noção dos processos que envolvem seu desenvolvimento: da ideia original à pesquisa, escrita, revisões, reescritas, pré-impressão, mais revisões, editoração, definição da capa, impressão final, distribuição e propaganda. Esse conjunto de aprendizados facilita o caminho de desenvolvimento de novos projetos.

\textbf{4. O que poderiam dizer como incentivo para que mais pesquisadores invistam na publicação de obras conscienciológicas?}

Uma Ciência não se desenvolve por si. A cada pesquisa, novas informações reiteram, avançam ou refutam informações pretéritas. É do esforço de cada um que cultivaremos, expandiremos, retribuiremos e colheremos frutos do trabalho conjunto.


\begin{pullquote}
``É do esforço de cada um que cultivaremos, expandiremos, retribuiremos e colheremos frutos do trabalho conjunto.''
\end{pullquote}

    
    \end{multicols}
\end{document}
